% M-x occur RET %%% RET

\DocumentMetadata{pdfstandard=A-2b}

\documentclass[12pt, a4paper, twoside]{memoir}
\usepackage[externalize=true, font=palatino, osf, math=fancy]{styles/style}

\pdfoutput=1
\pdfgentounicode=1

\hypersetup{%
  pdfa,
  pdftitle={Categorical Reconstruction Theory},
  pdfauthor={Tony Zorman},
  pdfdate={2025-04-01},
  pdfcopyright={Copyright (C) 2025, Tony Zorman},
  pdfsubject={Mathematics},
  pdfkeywords={Category theory, representation theory, Hopf algebras, monads, reconstruction theory},
  pdflicenseurl={https://creativecommons.org/licenses/by-sa/4.0/},
  pdfpubtype={thesis},
  pdflang={en},
  pdfmetalang={en},
}

\newcommand*{\submissionDate}{01.\ April 2025}% Used on title page and declaration.

%%% Colours for printing

\usepackage{xcolor}
\selectcolormodel{cmyk}
\definecolor{blue}{cmyk}{0.62, 0.49, 0, 0}
\definecolor{red}{cmyk}{0, 0.74, 0.79, 0.25}
\definecolor{green}{cmyk}{0, 0.72, 0.78, 0.31}

%%% Fixes for packages

% XXX: Probably better to add hooks as advised in [1], but for now
% redefining the internal command is easier for me.
% https://tex.stackexchange.com/questions/730148/cref-refers-to-lemmas-as-theorems
\makeatletter
\def\cref@thmoptarg[#1]#2#3#4{%
    \ifhmode\unskip\unskip\par\fi%
    \normalfont%
    \trivlist%
    \let\thmheadnl\relax%
    \let\thm@swap\@gobble%
    \thm@notefont{\fontseries\mddefault\upshape}%
    \thm@headpunct{.}% add period after heading
    \thm@headsep 5\p@ plus\p@ minus\p@\relax%
    \thm@space@setup%
    #2% style overrides
    \@topsep \thm@preskip               % used by thm head
    \@topsepadd \thm@postskip           % used by \@endparenv
    \def\@tempa{#3}\ifx\@empty\@tempa%
      \def\@tempa{\@oparg{\@begintheorem{#4}{}}[]}%
    \else%
      \refstepcounter[#1]{#3}%  <<< cleveref modification
      \@namedef{cref@#3@alias}{#1}% added
      \def\@tempa{\@oparg{\@begintheorem{#4}{\csname the#3\endcsname}}[]}%
    \fi%
    \@tempa}%
\makeatother

%%% Sidenotes
% style.sty has a `sidenotes' option, but the default sidenotes.sty file
% that CTAN has conflicts with the Bringhurst style of memoir. As such,
% we vendor sidenotes.sty to fix the incompatibilities, and manually
% redefine the sidenote-specific commands here.
\usepackage{sidenotes}
\usepackage{mparhack}
\renewcommand*{\marginfont}{\scriptsize}
\renewcommand{\note}[1]{%
  \sidenote{%
    \scriptsize
    \checkoddpage%
    \ifoddpage%
      \raggedright%
    \else
      \raggedleft%
    \fi
    {#1}%
  }%
}
\renewcommand{\scaption}[2][0cm]{%
  \marginnote{%
    \vspace{#1}%
    \caption{%
      \checkoddpage%
      \ifoddpage%
        \raggedright%
      \else
        \raggedleft%
      \fi
      {#2}%
    }%
  }
}

% Sidecaptions for figures get a small label, and an even slightly
% smaller body font.
\usepackage{caption}
\captionsetup{%
  style=base,%
  labelfont=footnotesize,%
  textfont=scriptsize,%
}

% Protrude equation numbers into the margin.
%
% Sources:
%   - https://tex.stackexchange.com/questions/258574/placing-the-equation-number-in-the-left-hand-margin
%   - https://www.overleaf.com/learn/latex/Page_size_and_margins
\makeatletter
\let\oldmaketag@@@\maketag@@@%
\def\oldtagform@#1{\oldmaketag@@@{(\ignorespaces#1\unskip\@@italiccorr)}}
\renewcommand{\eqref}[1]{\textup{\oldtagform@{\ref{#1}}}}
\newlength{\width@@}
\def\maketag@@@#1{\hbox{\hskip1sp\m@th\llap{%
      \normalfont#1%
      \settowidth{\width@@}{#1}%
      \checkoddpage\ifoddpage\hspace{-\the\width@@-\the\marginparsep}
      \else\hspace{\textwidth+\the\marginparsep+.1cm}\fi
    }}}
\makeatother

% Page separators for the introduction
\usepackage{froufrou}

%%% Additional maths macros

% No field for this category (as is the case in the style file), because
% things become much too overloaded otherwise.
\renewcommand*{\kvect}{\mathsf{vect}}
\renewcommand*{\kVect}{\mathsf{Vect}}

% Haven't gotten around to properly integrating these into style.sty, or
% they seem too specialised to do so.
\DeclareMathOperator{\Pic}{Pic}
\DeclareMathOperator{\QPiv}{QPiv}
\DeclareMathOperator{\Piv}{Piv}
\newcommand{\YangBaxter}{Yang--Baxter}
\newcommand*{\bbz}[1][B]{\ensuremath{{#1}^{\bullet}_0}\xspace}
\newcommand*{\bcz}[1][B]{\ensuremath{{#1}^{\circ}_0}\xspace}
\newcommand*{\bbt}[1][B]{\ensuremath{{#1}^{\bullet}_2}\xspace}
\newcommand*{\bct}[1][B]{\ensuremath{{#1}^{\circ}_2}\xspace}
\newcommand*{\duoidal}[1][D]{\ensuremath{(\cat{#1}, \circ, \bot, \bullet, 1)}\xspace}
\newcommand*{\duoidalb}[1][D]{\ensuremath{(\cat{#1}, \bullet, 1)}\xspace}
\newcommand*{\duoidalw}[1][D]{\ensuremath{(\cat{#1}, \circ, \bot)}\xspace}
\renewcommand*{\hom}[1]{\ensuremath{\lfloor#1\rfloor}}
\newcommand*{\cohom}[1]{\ensuremath{\lceil#1\rceil}}
\newcommand*{\rcomod}[1]{\ensuremath{\textnormal{comod-}#1}}
\newcommand*{\rComod}[1]{\ensuremath{\textnormal{Comod-}#1}}
\newcommand*{\psh}[1][C]{\widehat{\phantom{{}}\cat{#1}}} % presheaves
\newcommand*{\Cpsh}[1][C]{\widehat{\cat{#1}^{\op}}} % copresheaves
\newcommand*{\fin}{\mathrm{fin}}
\DeclareMathOperator{\Hom}{Hom}
\DeclareMathOperator{\tr}{tr}
\DeclareMathOperator{\GL}{GL}
\DeclareMathOperator{\End}{End}
\newcommand*{\Cocts}[1]{\ensuremath{#1\textnormal{-}\mathsf{Cocts}}}
\newcommand*{\Spiso}[1]{\mathsf{Sp}_{#1}}
\newcommand{\leftmod}[2]{#1\textnormal{-}\mathsf{Mod}_{#2}}
\newcommand{\modma}{\leftmod{A}{\mathcal{M}}}
\newcommand{\rightcomod}[2]{\mathsf{Comod}_{#1}{#2}}
\newcommand{\leftcomod}[2]{#1\textnormal{-}\mathsf{Comod}_{#2}}
\newcommand{\comodmc}{\leftcomod{C}{\mathcal{M}}}
\newcommand{\freeleftmod}[2]{#1\textnormal{-Free}_{#2}}
\newcommand{\forgetleftmod}[2]{#1\textnormal{-Forget}_{#2}}
\newcommand{\freeleftcomod}[2]{#1\textnormal{-Cofree}_{#2}}
\newcommand{\forgetleftcomod}[2]{#1\textnormal{-Coforget}_{#2}}
\newcommand*{\Ind}{\mathsf{Ind}}
\newcommand{\xtwoheadrightarrow}[2][]{%
  \xrightarrow[#1]{#2}\mathrel{\mkern-14mu}\rightarrow%
}
\newcommand{\biact}[2]{#1\textnormal{-}\mathsf{Biact}\textnormal{-}#2}
\newcommand{\biactC}[3][\cat{C}]{#2\textnormal{-}\mathsf{Biact}_{#1}\textnormal{-}#3}
\newcommand*{\bbcomod}[1][B]{\Tetramod[#1]}
\newcommand*{\bmodule}[1][B]{\tetramodfd[][#1][][]}
\newcommand*{\bbmodule}[1][B]{\Tetramod[][#1][][]}
\newcommand*{\EilenbergMoore}{\EM}
\newcommand*{\terminal}{\mathord{\heartsuit}}
\newcommand*{\blackdiamond}{\scriptscriptstyle\blacklozenge}
\NewDocumentCommand{\ZCat}{s o m o}{\buildCenter{#1}[Z][#2]{#3}[#4]}
\NewDocumentCommand{\ACat}{s o m o}{\buildCenter{#1}[A][#2]{#3}[#4]}
\NewDocumentCommand{\QCat}{s o m o}{\buildCenter{#1}[Q][#2]{#3}[#4]}
\NewDocumentCommand{\SZCat}{s o m o}{\buildCenter{#1}[SZ][#2]{#3}[#4]}
\newcommand*{\lbiDual}[1]{\lld{#1}}
\DeclareMathOperator{\Gr}{Gr}
\DeclareMathOperator{\Char}{Char}
\DeclareMathOperator{\Ad}{Ad}
\newcommand{\Malcev}{Mal'cev\xspace}
\DeclareMathOperator{\Isospace}{Iso}
\newcommand*{\modc}[2][\cat{C}]{\ensuremath{\mathrm{mod}_{#1}(#2)}}

%%% Additional environments

% The dedication: https://tex.stackexchange.com/questions/102830/dedication-page-in-article-class
\newenvironment{dedication}{%
  \cleardoublepage%      we want a new page
  \thispagestyle{empty}% no header and footer
  \vspace*{\stretch{1}}% some space at the top
  \ttfamily%             a monospaced font for specialness
  \raggedleft%           flush to the right margin
}{%
  \par%                  end the paragraph
  \vspace{\stretch{3}}%  space at bottom is three times that at the top
}

\theoremstyle{plain}
\newtheorem{theoremalph}{Theorem}
\renewcommand*{\thetheoremalph}{\Alph{theoremalph}}

% Create a fake theorem environment for when they appear in the
% introduction, so one doesn't have to write Theorem (Theorem x.y) or
% other such constructs, but directly creates a theorem with a link to
% the "real" one.
%
% Original source: https://tex.stackexchange.com/questions/422/how-do-i-repeat-a-theorem-number
\makeatletter
\newtheorem*{rep@theorem}{\rep@title}
\newenvironment{reptheorem}[1]{%
  \def\rep@title{\cref{#1}}\begin{rep@theorem}%
}{%
  \end{rep@theorem}%
}
\makeatother

% There's a joke in here somewhere. Anyways, use realhats, but only with
% \Hat—the original \hat should stay the way it is.
\let\oldHat\hat%
\usepackage{realhats}
\let\Hat\hat%
\let\hat\oldHat%

%%% Changes to memoir

% Chapter style; memoir prominently features a "Bringhurst" style, but
% neglects to add the large chapter numbers into the margin, as the
% famous book does. Fix this.
%
% Source: https://tex.stackexchange.com/questions/88895/bringhurst-chapter-style-in-memoir
\usepackage{bringhurst}

% Table of contents
\renewcommand{\cftchapterfont}{\scshape} % Small-caps for chapters
\renewcommand{\cftchapterpagefont}{}     % Nothing for sections
\setcounter{tocdepth}{4}

% Floats that are alone on a page get an empty style, and otherwise we
% want a minimal-esque ruled one.
\mergepagefloatstyle{dissstyle}{ruled}{empty}
\pagestyle{dissstyle}

% Epigraphs
\epigraphfontsize{\footnotesize}

%%%% Finalise the page layout.

\usepackage{flafter}% Figures should appear after they are defined.

% This page layout is a little wider than the default one memoir picks
% for a4paper, which means we can increase the font size a bit. With
% 12pt Palatino, this results in about 72 characters per full line,
% which is on the higher side of what I would prefer, but either
% increasing the font size or decreasing the margins even more would
% impede too much upon largeish diagrams or formulas, of which this
% theses has enough of.
\isopage%

% Do not give sidenotes the full width of the margins, as that looks a
% bit weird in my opinion. This results in a maximum of about 21
% characters—small, but manageable.
\setmarginnotes{17pt}{80pt}{\onelineskip}

% Instead of
%
%     \clubpenalty=10000
%     \widowpenalty=10000
%     \raggedbottom%
%
% we use a larger topskip, together with the \sloppybottom below.
% Empirically, this produces slightly nicer looking pages.
\setlength{\topskip}{1.6\topskip}

% This must run after all changes have been applied.
\checkandfixthelayout%

\sloppybottom%

%%% Local Variables:
%%% mode: LaTeX
%%% TeX-master: "../main"
%%% End:
